\documentclass[12pt,a4paper]{article}

% ================= PACOTES =================
\usepackage[utf8]{inputenc}
\usepackage[T1]{fontenc}
\usepackage{amsmath,amssymb,amsfonts}
\usepackage{geometry}
\usepackage{graphicx}
\usepackage{float}
\usepackage{hyperref}
\usepackage{xcolor}
\usepackage{booktabs}
\usepackage{enumitem}
\usepackage{fancyhdr}
\usepackage{titlesec}
\usepackage{caption}

\geometry{margin=2.5cm}

\hypersetup{
    colorlinks=true,
    linkcolor=blue,
    citecolor=blue,
    urlcolor=blue
}

\pagestyle{fancy}
\fancyhf{}
\rhead{APC \#2 -- Cálculo Numérico Aplicado}
\lhead{\thepage}
\renewcommand{\headrulewidth}{0.4pt}

\titleformat{\section}{\normalfont\Large\bfseries}{\thesection}{1em}{}
\titleformat{\subsection}{\normalfont\large\bfseries}{\thesubsection}{1em}{}

% ================= DADOS DO ALUNO =================
\newcommand{\nomealuno}{NOME DO ALUNO}
\newcommand{\matriculaaluno}{MATRÍCULA}
\newcommand{\datadoc}{\today}

\title{
    \vspace{-1.5cm}
    \textbf{Atividade para Casa \#2 (APC2)} \\
    \large Sistemas Dinâmicos com Múltiplos Graus de Liberdade \\
    \large e o Polinômio Característico
}
\author{\nomealuno \\ Matrícula: \matriculaaluno \\[0.3cm]
    Cálculo Numérico Aplicado \\
    Prof. Doutor Rafael Gabler Gontijo}
\date{\datadoc}

\begin{document}

\maketitle

\begin{abstract}
\noindent Este relatório apresenta o desenvolvimento analítico solicitado na Atividade
para Casa \#2 (APC2), cujo objetivo é estabelecer a conexão entre a modelagem física de
sistemas massa-mola-amortecedor com $N$ graus de liberdade e a estrutura matemática do
problema de autovalores generalizado associado, culminando na obtenção do polinômio
característico do sistema. Os resultados aqui obtidos servirão de base para a
implementação computacional do método de Bairstow no Programa para Casa \#2 (PPC2).
\end{abstract}

\tableofcontents
\newpage

% =====================================================
\section{Introdução}
% =====================================================

% TODO: Escreva uma breve introdução contextualizando o problema:
% - sistemas massa-mola-amortecedor com N graus de liberdade
% - importância de polinômios característicos na engenharia
% - conexão com o método de Bairstow (a ser implementado no PPC2)

Texto introdutório aqui...

% =====================================================
\section{Formulação Geral do Problema}
% =====================================================

A equação do movimento de um sistema massa-mola-amortecedor com $N$ graus de
liberdade, na ausência de forças externas, é escrita na forma matricial como:

\begin{equation}
    M\ddot{\mathbf{x}}(t) + C\dot{\mathbf{x}}(t) + K\mathbf{x}(t) = \mathbf{0},
    \label{eq:eom}
\end{equation}

\noindent em que $M, C, K \in \mathbb{R}^{N\times N}$ são, respectivamente, as
matrizes de massa, amortecimento e rigidez, e $\mathbf{x}(t) \in \mathbb{R}^N$ é o
vetor de deslocamentos generalizados.

% =====================================================
\section{Item 1 -- Obtenção do Polinômio Característico de Ordem $2N$}
% =====================================================

\subsection{Proposta de solução}

Assumindo uma solução do tipo:

\begin{equation}
    \mathbf{x}(t) = \boldsymbol{\phi}\, e^{\lambda t},
    \label{eq:ansatz}
\end{equation}

\noindent em que $\boldsymbol{\phi} \in \mathbb{C}^N$ é um vetor constante
(o modo de vibração) e $\lambda \in \mathbb{C}$ é o autovalor a ser determinado.

\subsection{Substituição na equação do movimento}

% TODO: desenvolva aqui, passo a passo, a substituição de (2) em (1):
% - calcule dx/dt e d^2x/dt^2
% - substitua na equação do movimento
% - cancele o termo exponencial (justifique por que isso é válido)

Derivando a Eq.~\eqref{eq:ansatz} em relação ao tempo:

\begin{equation}
    \dot{\mathbf{x}}(t) = \lambda\,\boldsymbol{\phi}\, e^{\lambda t}, \qquad
    \ddot{\mathbf{x}}(t) = \lambda^2\,\boldsymbol{\phi}\, e^{\lambda t}.
    \label{eq:derivadas}
\end{equation}

Substituindo \eqref{eq:ansatz} e \eqref{eq:derivadas} na equação do movimento
\eqref{eq:eom}:

\begin{equation}
    \left(\lambda^2 M + \lambda C + K\right)\boldsymbol{\phi}\, e^{\lambda t}
    = \mathbf{0}.
    \label{eq:subst}
\end{equation}

Como $e^{\lambda t} \neq 0$ para todo $t$, obtemos o \textbf{problema de
autovalores generalizado}:

\begin{equation}
    \boxed{\left(\lambda^2 M + \lambda C + K\right)\boldsymbol{\phi} = \mathbf{0}.}
    \label{eq:autovalor_generalizado}
\end{equation}

\subsection{Condição para solução não trivial}

% TODO: justifique por que phi = 0 não interessa fisicamente,
% e por que isso implica que a matriz precisa ser singular

Para que exista solução $\boldsymbol{\phi} \neq \mathbf{0}$ (solução não trivial,
única de interesse físico), a matriz $A(\lambda) = \lambda^2 M + \lambda C + K$
deve ser singular, isto é:

\begin{equation}
    \det\left(\lambda^2 M + \lambda C + K\right) = 0.
    \label{eq:det_zero}
\end{equation}

\subsection{Grau do polinômio característico}

% TODO: demonstre, usando a definição de determinante (regra de Leibniz)
% ou a estrutura da matriz A(lambda), que o grau máximo do polinômio
% resultante é 2N, e que o coeficiente do termo de maior grau é det(M)

Cada entrada da matriz $A(\lambda) = \lambda^2 M + \lambda C + K$ é um polinômio
de grau, no máximo, $2$ em $\lambda$. Como o determinante de uma matriz
$N \times N$ é uma soma de produtos de $N$ entradas (uma por linha e por coluna),
cada termo dessa soma tem grau no máximo $2N$. O termo de maior grau resulta do
produto dos elementos da diagonal principal, o que fornece o coeficiente
líder $\det(M)$. Assim, o polinômio característico é dado por:

\begin{equation}
    P(\lambda) = \det\left(\lambda^2 M + \lambda C + K\right)
    = a_{2N}\lambda^{2N} + a_{2N-1}\lambda^{2N-1} + \cdots + a_1\lambda + a_0,
    \label{eq:poli_geral}
\end{equation}

\noindent com $a_{2N} = \det(M) \neq 0$ (assumindo $M$ não singular), garantindo
que $P(\lambda)$ tenha grau exatamente $2N$.

% =====================================================
\section{Item 2 -- Significado Físico das Raízes}
% =====================================================

% TODO: discuta o significado físico das 2N raízes de P(lambda):
% - raízes complexas conjugadas -> modos oscilatórios amortecidos
% - parte real negativa -> estabilidade / decaimento
% - parte real positiva -> instabilidade
% - raízes reais puras -> comportamento superamortecido / não oscilatório
% - relação entre |Im(lambda)| e a frequência natural amortecida
% - relação entre Re(lambda) e a taxa de decaimento / fator de amortecimento

Discussão física aqui...

% =====================================================
\section{Item 3 -- Sistema Numérico Proposto}
% =====================================================

% TODO: escolha N (sugestão: N = 2 ou N = 3) e proponha valores numéricos
% para massas m_i, coeficientes de amortecimento c_i e rigidezes k_i.
% Desenhe (ou descreva) o sistema massa-mola-amortecedor correspondente.

\subsection{Descrição do sistema}

Considere um sistema com $N = \_\_$ graus de liberdade, composto pelas massas:

\begin{equation}
    m_1 = \_\_, \quad m_2 = \_\_, \quad \ldots
\end{equation}

\noindent conectadas pelas rigidezes:

\begin{equation}
    k_1 = \_\_, \quad k_2 = \_\_, \quad \ldots
\end{equation}

\noindent e pelos coeficientes de amortecimento:

\begin{equation}
    c_1 = \_\_, \quad c_2 = \_\_, \quad \ldots
\end{equation}

% Sugestão: inclua aqui uma figura esquemática do sistema (TikZ) se desejar.
% \begin{figure}[H]
%     \centering
%     % \input{seu_diagrama_tikz.tex}
%     \caption{Esquema do sistema massa-mola-amortecedor com $N$ graus de liberdade.}
%     \label{fig:sistema}
% \end{figure}

% =====================================================
\section{Item 4 -- Construção Explícita do Sistema e do Polinômio Característico}
% =====================================================

% TODO: monte explicitamente as matrizes M, C, K para os valores escolhidos

\subsection{Matrizes do sistema}

As matrizes de massa, amortecimento e rigidez para o sistema proposto são:

\begin{equation}
    M =
    \begin{bmatrix}
        \_ & \_ \\
        \_ & \_
    \end{bmatrix}, \qquad
    C =
    \begin{bmatrix}
        \_ & \_ \\
        \_ & \_
    \end{bmatrix}, \qquad
    K =
    \begin{bmatrix}
        \_ & \_ \\
        \_ & \_
    \end{bmatrix}.
    \label{eq:matrizes}
\end{equation}

\subsection{Determinante e polinômio característico}

% TODO: calcule explicitamente det(lambda^2 M + lambda C + K) passo a passo
% para o seu caso (N=2 -> polinômio de grau 4; N=3 -> polinômio de grau 6)

\begin{equation}
    \det\left(\lambda^2 M + \lambda C + K\right) = \_\_\_\_
\end{equation}

Desenvolvendo os produtos e agrupando os termos por potência de $\lambda$:

\begin{equation}
    P(\lambda) = \_\_ \lambda^{2N} + \_\_ \lambda^{2N-1} + \cdots + \_\_ \lambda + \_\_
\end{equation}

% =====================================================
\section{Item 5 -- Polinômio na Forma Padrão}
% =====================================================

% TODO: organize o polinômio final na forma padrão pedida, explicitando
% numericamente cada coeficiente a_i, pronto para ser usado no PPC2.

O polinômio característico obtido pode ser reescrito na forma padrão:

\begin{equation}
    P(\lambda) = a_{2N}\lambda^{2N} + a_{2N-1}\lambda^{2N-1} + \cdots
    + a_1\lambda + a_0,
\end{equation}

\noindent com os coeficientes numéricos:

\begin{table}[H]
    \centering
    \caption{Coeficientes do polinômio característico.}
    \label{tab:coeficientes}
    \begin{tabular}{cc}
        \toprule
        Coeficiente & Valor \\
        \midrule
        $a_{2N}$ & \_\_ \\
        $a_{2N-1}$ & \_\_ \\
        $\vdots$ & $\vdots$ \\
        $a_1$ & \_\_ \\
        $a_0$ & \_\_ \\
        \bottomrule
    \end{tabular}
\end{table}

Esses coeficientes serão utilizados diretamente na implementação do método de
Bairstow no Programa para Casa \#2 (PPC2), para a determinação numérica das
raízes $\lambda_i$ e, consequentemente, dos modos dinâmicos do sistema.

% =====================================================
\section{Conclusão}
% =====================================================

% TODO: sintetize os principais resultados obtidos e a conexão com o PPC2

Conclusão aqui...

% =====================================================
\section*{Referências Bibliográficas}
% =====================================================

\begin{enumerate}[label={[\arabic*]}]
    \item S. C. Chapra, R. P. Canale. \textit{Métodos Numéricos para Engenharia}.
    McGrawHill, 5\textsuperscript{a} edição, 2008.
    \item R. G. Gontijo. \textit{Notas de aula do curso de Cálculo Numérico
    Aplicado}, 2026.
\end{enumerate}

\end{document}
